\begin{afterword}
\summaryhead

The post offers a trick for questions that seem impossible to answer: instead of asking ``Why do I have free will?'', ask ``Why do I think I have free will?'' The second question always has an answer, because the belief is a real event with causes, whether or not free will exists or even makes sense. The same substitution works for time, consciousness and the existence of reality.

The author illustrates with socks. Tracing the belief ``I am wearing socks'' back through the visual cortex, the retina and reflected light ends at the socks themselves. A mirage of a lake, by contrast, is explained without any lake: the belief is explained away. Saying ``I believe in free will because I have free will'' is not enough; each step of the causal chain must be explained in terms that are not themselves confusing. Where, the post asks, is the free-will sensor? Questions of cognitive science, it concludes, are at least solvable.

\responsehead

The trick is clear and useful, and the post states it well. A belief is an event, and asking what caused it gives the inquiry a definite starting point even when the original question is confused. The socks and the mirage show the two ways the inquiry can end.

The method has a history that the post does not mention. Hume proposed a close relative in 1748: when a term is suspected of having no meaning, ``we need but enquire, from what impression is that supposed idea derived?'' For free will, a cognitive account of how the feeling of willing arises had been published six years before the post. For consciousness, the post's pair of questions is close to what David Chalmers later called the meta-problem, the question of why we think there is a problem of consciousness, and Chalmers finds versions of the strategy in Hobbes, Hume, Spinoza and Kant.

The guarantee covers less than it may seem to. It is a guarantee that the question about the belief has an answer. Whether that answer settles the original question depends on where the causal chain ends, and on being able to recognize the end. For socks this is easy. For consciousness, one of the post's own examples, it is exactly what is disputed: illusionists, who hold that consciousness is not what it seems to us, typically hold that explaining the belief dissolves the problem, and Chalmers holds that it ``does not suffice to solve or dissolve the problem of consciousness.'' The post's closing claim that questions of cognitive science are ``solvable'' is true of the substitute question. It does not show that the original questions are thereby solved.

In short: a clear statement of an old and useful method, whose guarantee applies to the question about the belief and does not by itself settle the question the reader started with.
\end{afterword}
